% Part: counterfactuals
% Chapter: introduction
% Section: counterfactuals

\documentclass[../../../include/open-logic-section]{subfiles}

\begin{document}

\olfileid[zh]{cnt}{int}{cnt}

\olsection{反事实条件句}

英语中一种非常常见且重要的「if \dots then \dots」构式，是利用系动词
\emph{to be}的过去虚拟式构成的：「if it were the case that \dots then it would be the
case that \dots」。由于这样的条件句的前件通常为假，即与事实相悖，故称之为\emph{反事实条件句}（并且因为它们使用了系动词 \emph{to be}的虚拟式，也称作\emph{虚拟条件句}）。它们与采取「if it is the case that \dots then it is the case that \dots」形式的\emph{直陈}条件句相区别。反事实与直陈条件句在真值条件上有所不同。考虑 Adams 的著名例子：
\begin{quote}
  如果 Oswald 没有杀 Kennedy，那么是别人杀的。

  如果 Oswald 没有杀 Kennedy，那么别人也会杀他。
\end{quote}
第一个是直陈的，第二个是反事实的。第一个显然为真：我们知道总统 John F. Kennedy 是被\emph{某个人}杀害的，而如果这个人不是（与 Warren 报告相反）Lee Harvey Oswald，那么就是别人杀死了 Kennedy。第二个说的则是不同的事情。它主张，如果 Oswald 没有杀 Kennedy，即如果达拉斯枪击事件被避免或未能成功，那么历史后续就会以这样一种方式展开：另一次暗杀将会成功。要使它为真，就必须存在着强大的势力密谋以确保 JFK 之死（正如许多 JFK 阴谋论者所相信的那样）。

关于\emph{直陈}条件句是否能被实质条件句正确地刻画，是一个仍在进行的争论，尤其是关于实质条件句的悖论能否以一种与之相容、从而为英语直陈条件句给出真值条件的方式被「解释」。相比之下，反事实条件句不能被实质条件句正确符号化则是不争的事实。这一点是明确的，因为尽管反事实条件句的前件通常为假，但并非所有前件为假的反事实条件句都为真——例如，如果你相信 Warren 报告，且并不存在暗杀 JFK 的阴谋，那么 Adams 的反事实条件句就是这样一个反例。

反事实条件句在因果推理中扮演重要角色：使用反事实条件句的一个典型例子就是表达因果关系。例如，划火柴会使其点燃，你可以通过说「如果这根火柴被划了，它就会亮起来」来表述这一点。实质条件句以及一般而言的直陈条件句不能用来表述这一点：「火柴被划了 $\lif$ 火柴亮了」在火柴从未被划的情况下为真，无论如果被划了会发生什么都是如此。更糟的是，「火柴被划了 $\lif$ 火柴变成一束花」在火柴从未被划的情况下也为真，但火柴如果被划了当然不会变成一束花。

关于反事实条件句的正确逻辑究竟是什么，仍在争论之中。\zhFullName{Robert Stalnaker}{罗伯特·斯达内克}{斯达内克}和\zhFullName{David K. Lewis}{大卫·刘易斯}{刘易斯}给出了一个有影响力的反事实条件句分析。按照他们的看法，一个反事实条件句「如果~$S$ 是这种情况，那么~$T$ 就是这种情况」为真，当且仅当 $T$ 在与现实世界最接近且~$S$ 为真的那个反事实情形（「可能世界」）中为真。这被称为「本体性」分析，因为它涉及可能世界的本体论。其他分析则利用条件概率或信念修正理论。目前涌现了各种不同的反事实条件句逻辑。甚至不存在唯一一种\zhFullName{David K. Lewis}{大卫·刘易斯}{刘易斯}--\zhFullName{Robert Stalnaker}{罗伯特·斯达内克}{斯达内克}反事实条件句逻辑：尽管\zhFullName{Robert Stalnaker}{罗伯特·斯达内克}{斯达内克}和\zhFullName{David K. Lewis}{大卫·刘易斯}{刘易斯}参照最接近的可能世界提出了思路相近的论述，但他们所做的假设导致了不同的有效推理。

\end{document}
